\paragraph{Accuracy of the learned value head.}
To separate the quality of the learned Actor from the improvement supplied by
the final stochastic evaluator, Tables~\ref{tab:direct_value_head_three_pdes}
and~\ref{tab:joint_nonseparable_value_head} evaluate the selected Actor's
deterministic value head directly, without an FK rollout or martingale
correction.  On the matched 960-state whole-function banks in
Table~\ref{tab:direct_value_head_three_pdes}, the direct head ranks first in
$6/15$ PDE--dimension cells and second in another $8/15$, placing in the top
two in $14/15$ cells.  It wins $35/45$ matched pairwise comparisons against
the independently trained baselines: $7/15$ against DPI, $13/15$ against
PINN-HTE, and $15/15$ against SDGD-PINN.  Thus, the value head is already a
competitive deterministic whole-function approximator, although it is not
uniformly best.  The separate 32-state single-point diagnostic in
Table~\ref{tab:joint_nonseparable_value_head} gives a complementary result:
for the primary RMSE, the direct head outperforms FH-MLP in all five dimensions
and every numerically valid Deep BSDE result ($3/3$), while the complete
FK--MCV estimator remains more accurate in all five dimensions; Deep BSDE has
numerical failures at $d=200$ and $d=400$.  RelMAE gives the same broad
ordering, except that Deep BSDE is slightly better than the direct head at
$d=20$.  Because the two tables use different state banks, baseline sets, and
replication protocols, their absolute errors should not be compared directly.
Together, they show that the learned Actor itself carries substantial solution
accuracy, while the final FK--MCV evaluation provides a further, consistent
accuracy gain.
